%! The Tangent Function — Unit 3, Lesson 3.8 — Homework
\documentclass[10pt]{article}
\usepackage{precalculus-article}
\usepackage{precalculus-boxes}
\newcommand{\TallMath}[1]{$\displaystyle #1\rule[-1.4em]{0pt}{3.2em}$}

\begin{document}

\pageheader{Unit 3, Lesson 3.8}{Homework}
\namedateperiod

\vspace{0.08in}

\begin{notesbox}{Problems 1--6}

\textbf{1.}\quad Use the unit circle to evaluate each expression.  If the value
is undefined, write ``undefined.''

\vspace{4pt}
\renewcommand{\arraystretch}{1.5}
\begin{tabular}{@{} p{0.28\linewidth} p{0.28\linewidth} p{0.28\linewidth} @{}}
(a) $\tan(0) = $ \blank{1.6cm} &
(b) $\tan\!\left(\dfrac{\pi}{4}\right) = $ \blank{1.4cm} &
(c) $\tan\!\left(\dfrac{\pi}{2}\right) = $ \blank{1.4cm} \\[12pt]
(d) $\tan(\pi) = $ \blank{1.6cm} &
(e) $\tan\!\left(\dfrac{3\pi}{4}\right) = $ \blank{1.4cm} &
(f) $\tan\!\left(-\dfrac{\pi}{4}\right) = $ \blank{1.4cm} \\
\end{tabular}

\vspace{0.12in}
\textbf{2.}\quad For each function, state the \textbf{period} and the
location of \textbf{all} vertical asymptotes (write a general formula
involving integer $k$).

\vspace{4pt}
(a)\quad $g(\theta) = \tan\!\left(\dfrac{\theta}{3}\right)$

\vspace{4pt}
\hspace{1em} Period $= $ \blank{3cm} \hfill
Asymptotes: $\theta = $ \blank{5cm}

\vspace{0.10in}
(b)\quad $h(\theta) = \tan(4\theta)$

\vspace{4pt}
\hspace{1em} Period $= $ \blank{3cm} \hfill
Asymptotes: $\theta = $ \blank{5cm}

\vspace{0.10in}
(c)\quad $p(\theta) = -3\tan(\theta) + 2$

\vspace{4pt}
\hspace{1em} Period $= $ \blank{3cm} \hfill
Asymptotes: $\theta = $ \blank{5cm}

\vspace{0.12in}
\textbf{3.}\quad A function is given by
$f(\theta) = \tan\!\left(\theta - \dfrac{\pi}{6}\right)$.

\vspace{4pt}
(a)\quad What is the phase shift?
\writeline

\vspace{4pt}
(b)\quad Find the vertical asymptote closest to $\theta = 0$ with $\theta > 0$.
Show your work.
\writelines{3}

\vspace{0.12in}
\textbf{4.}\quad Describe in words how the graph of $g(\theta) = 3\tan(2\theta)$
differs from the graph of $f(\theta) = \tan(\theta)$.  Address the period and
vertical scale.
\writelines{3}

\vspace{0.12in}
\textbf{5.}\quad \textbf{(AP-Style MC)}\quad Which of the following is a vertical
asymptote of $g(\theta) = \tan\!\left(\dfrac{\theta}{2} + \dfrac{\pi}{4}\right)$?

\begin{enumerate}[label=(\Alph*), leftmargin=2em, itemsep=3pt]
  \item $\theta = \dfrac{\pi}{4}$
  \item $\theta = \dfrac{\pi}{2}$
  \item $\theta = \pi$
  \item $\theta = \dfrac{3\pi}{2}$
  \item $\theta = 2\pi$
\end{enumerate}

Answer: \underline{\hspace{1.5cm}}

\vspace{0.12in}
\textbf{6.}\quad The table below shows values of a function $q$.

\vspace{4pt}
\begin{center}
\renewcommand{\arraystretch}{1.4}
\begin{tabular}{|c|c|c|c|c|c|}
\hline
$\theta$ & $-\pi/4$ & $0$ & $\pi/4$ & $\pi/2$ & $3\pi/4$ \\
\hline
$q(\theta)$ & $-1$ & $0$ & $1$ & undef. & $-1$ \\
\hline
\end{tabular}
\end{center}

\vspace{4pt}
(a)\quad Is $q$ consistent with $q(\theta) = \tan\theta$?  Explain.
\writelines{2}

(b)\quad What is the period of $q$ suggested by the table?
\writeline

\end{notesbox}

\vspace{0.06in}

\begin{extensionbox}
\textbf{Extension (Optional):}  Show algebraically that $\tan(\theta + \pi) = \tan\theta$
for all $\theta$ in the domain.  Start from
$\tan(\theta+\pi) = \dfrac{\sin(\theta+\pi)}{\cos(\theta+\pi)}$
and use the angle addition identities
$\sin(\theta+\pi) = -\sin\theta$ and $\cos(\theta+\pi)=-\cos\theta$.
\writelines{4}
\end{extensionbox}

\vspace{0.06in}

\begin{tcolorbox}[colback=navylight, colframe=navy, arc=2mm,
    left=3mm, right=3mm, top=2mm, bottom=2mm]
\small\textbf{Preview --- Lesson 3.9: Inverse Trigonometric Functions}\\
In Lesson 3.9 we will ask: if we know a sine, cosine, or tangent value, can we
find the angle?  To answer that question we restrict the domain of each
trigonometric function to make it invertible, then define $\arcsin$, $\arccos$,
and $\arctan$.  The asymptote behavior you studied today --- $\tan\theta\to\pm\infty$
as $\theta\to\pm\pi/2$ --- is precisely what tells us the range of $\arctan$.
\end{tcolorbox}

\end{document}
